\subsection{启动与访问}
从仓库根目录执行以下命令。启动时依次准备依赖、Raft 节点、事务服务和控制台。控制台入口为 \url{http://127.0.0.1:5173}，事务端口以 Compose 输出为准。浏览器接入同一个 Observer，通过对应页面读取节点、事务和归档数据。

\begin{lstlisting}
pnpm install --frozen-lockfile
pnpm demo
docker compose -f compose.transaction.yaml up -d --build transaction
docker compose -f compose.transaction.yaml port transaction 8080
pnpm console
\end{lstlisting}

在另一个终端可使用 \texttt{pnpm verify} 对归档历史生成 V01--V03 判定。也可指定来源历史，回放一次键值调用序列：

\begin{lstlisting}
pnpm verify V01 2026-11-07T12-03-08.334Z--8c1d3503
\end{lstlisting}

\subsection{运行记录索引}
表~\ref{tab:records} 的每行对应仓库 \path{evidence/} 下按编号和 Run ID 命名的目录。\path{metadata.json} 给出运行参数、记录时间和对应代码版本，\path{events.jsonl} 保存事件与调用历史，\path{result.json} 汇总判定及其依据。表中时间使用 UTC。

\begin{table}[htbp]
\centering
\footnotesize
\caption{正文引用的运行记录}
\label{tab:records}
\begin{tabularx}{\linewidth}{@{}l l X@{}}
\toprule
编号 & Run ID & 对应内容 \\
\midrule
SMOKE & \texttt{2026-11-07T12-02-33.941Z--13aad51c} & 五节点部署及键值读写 \\
R01 & \texttt{2026-11-06T10-57-02.353Z--76c65d01} & Leader 切换、已确认值与原节点重启 \\
R02 & \texttt{2026-11-06T10-57-06.906Z--77080c4c} & 分区两侧返回与前缀收敛 \\
R03 & \texttt{2026-11-07T12-03-08.334Z--8c1d3503} & 并发客户端调用历史 \\
E3B & \texttt{2026-11-07T10-50-57.016Z--9be3bb1d} & 日志压缩、快照安装与重启 \\
T01 & \texttt{2026-11-08T14-15-50.333Z--f198395b} & 提交边界、历史读取与检查点恢复 \\
T02 & \texttt{2026-11-07T13-41-40.891Z--561ddfc6} & 值班约束对应的事务历史 \\
T03 & \texttt{2026-11-08T14-17-06.819Z--58518223} & CLR 重做与继续撤销 \\
V01 & \texttt{2026-11-07T14-45-07.744Z--aa3b0215} & R03 的合法顺序 \\
V02 & \texttt{2026-11-07T14-45-07.833Z--4a533f0f} & T02 的版本依赖与环 \\
V03 & \texttt{2026-11-07T14-45-07.911Z--30b03c6e} & T02 的值班人数判定 \\
\bottomrule
\end{tabularx}
\end{table}
